\section{Writing}
\label{sec:writing}

\subsection{The editor}

The editing surface is Monaco --- the engine behind VS Code --- with a LaTeX
grammar written for this app. Commands, env names, maths and comments
are all coloured differently, \verb|\begin|/\verb|\end| blocks fold, and
\kbd{Cmd+/} comments the current line.

The editor's colours come from the same palette as the rest of the interface,
so switching between the dark and light themes in \ui{Settings} restyles the
editor along with everything else, rather than leaving a dark rectangle in a
light window.

\subsection{Snippets}

Start typing a construct's name and the completion list appears; press
\kbd{Tab} to move between the placeholders it leaves behind. Typing
\texttt{figure} and accepting the suggestion gives you:

\begin{lstlisting}
\begin{figure}[htbp]
	\centering
	\includegraphics[width=0.8\textwidth]{path}
	\caption{caption}
	\label{fig:label}
\end{figure}
\end{lstlisting}

with the cursor on \texttt{0.8\textbackslash textwidth} and \kbd{Tab} stepping
through the rest. The full set covers \texttt{begin}, \texttt{frac},
\texttt{sqrt}, \texttt{sum}, \texttt{int}, \texttt{lim}, \texttt{prod},
\texttt{section}, \texttt{subsection}, \texttt{itemize}, \texttt{enumerate},
\texttt{figure}, \texttt{table}, \texttt{align}, \texttt{matrix},
\texttt{cases}, \texttt{cite}, \texttt{ref}, and the text-formatting commands.

\subsection{The component library}

Snippets are skeletons; components are filled-in blocks. Open \ui{Components}
in the sidebar, pick one, complete the form, and press \ui{Insert} to place the
generated \LaTeX{} at your cursor.

\begin{table}[htbp]
  \centering
  \caption{What the component library offers}
  \label{tab:components}
  \begin{tabular}{ll}
    \toprule
    \textbf{Category} & \textbf{Examples} \\
    \midrule
    Mathematics        & numbered and aligned equations, matrices, cases, theorems \\
    Tables             & basic and publication tables with \texttt{booktabs} rules \\
    Figures            & single figures, side-by-side subfigures, Ti\emph{k}Z diagrams \\
    Lists              & itemised and enumerated lists \\
    Structure          & sections, abstracts \\
    Code               & source listings, algorithm and pseudocode blocks \\
    Bibliography       & article and book Bib\TeX{} entries \\
    \bottomrule
  \end{tabular}
\end{table}

Use the library when you want a correct starting point for something fiddly ---
a \texttt{booktabs} table with the rules in the right places, or a subfigure
pair with aligned captions --- and snippets when you already know the shape and
just want the keystrokes saved.

\subsection{Moving around a document}

The \ui{Outline} panel lists every heading in the open file and jumps to it on
click. It reads \verb|\part| through \verb|\subparagraph|, understands starred
forms like \verb|\section*|, and ignores anything commented out --- so a
heading you have temporarily disabled with a \texttt{\%} does not clutter it.

For longer jumps, the command palette (\kbd{Cmd+Shift+P}) exposes go-to-line,
find, find-and-replace, comment toggling, line duplication and movement, and
the editor zoom controls, and the full shortcut list --- search it for
\emph{keyboard shortcuts}. \kbd{Cmd+K} is not the way in: that key opens the
inline assistant, Section~\ref{sec:inline}.

\subsection{Reading errors}

When a run fails, the \ui{Issues} panel lists what went wrong: errors first,
then warnings, then overfull and underfull boxes. Each row gives the message,
the offending source line where the log provides it, and the file and line
number.

Click a row and the editor opens that file and puts your cursor on that line.

The distinction is worth respecting. An \emph{error} stopped the compiler and
must be fixed. A \emph{warning} --- an undefined reference, a missing citation
--- produced a PDF anyway, but usually one with \texttt{??} in it. A
\emph{badbox} is purely typographic: a line \LaTeX{} could not break neatly.
Badboxes are not clickable, because there is generally nothing at that line to
fix; they are a signal to reword a sentence or let it go.

Remember Section~\ref{sec:passes}: a fresh project reports undefined references
on its first run and resolves them on the second. Do not go hunting for a bug
that a second press of \ui{Typeset} will fix.
